% ============================ INTRODUCCIÓN ============================
\chapter{Introducción}\label{cap:introduccion}

\section{Motivación}
Este documento recoge el montaje y la puesta a prueba de un laboratorio SOC
(Security Operations Center) casero, desplegado íntegramente con
máquinas virtuales sobre un portátil con Ubuntu. El objetivo es reproducir a
pequeña escala el trabajo de un analista de nivel 1: recoger telemetría de los
equipos, generar ataques controlados y comprobar si se detectan, cómo se
detectan y qué se puede hacer como respuesta.

\pendiente{Breve explicación de las motivaciones del proyecto}

\section{Objetivos}
\begin{itemize}[leftmargin=*]
  \item Desplegar un SIEM/XDR funcional (Wazuh) en una red aislada.
  \item Segmentar la red con un cortafuegos (OPNsense) y añadir un IDS de red (Suricata).
  \item Monitorizar equipos víctima Windows y Linux con agentes y Sysmon.
  \item Ejecutar escenarios de ataque desde una red de atacante separada (Kali)
        y documentar detección, análisis y respuesta, mapeados a MITRE ATT\&CK.
\end{itemize}

\section{Alcance y limitaciones}
\begin{itemize}[leftmargin=*]
  \item Entorno de laboratorio sin tráfico real de usuarios: el volumen de
        eventos y falsos positivos no es representativo de producción.
  \item Recursos limitados del host (RAM), lo que condiciona cuántas VM pueden
        estar encendidas a la vez.
  \item Todos los ataques se ejecutan contra máquinas propias dentro de redes aisladas.
\end{itemize}

\begin{nota}[Hardware del equipo anfitrión]
\begin{table}[H]\centering\small
\begin{tabularx}{\linewidth}{>{\bfseries\color{socazul}}l X}
  \toprule
  Procesador     & Intel Core i7-8750H @ 2.20\,GHz (6 núcleos / 12 hilos) \\
  Memoria RAM    & 32\,GB DDR4 en dos módulos de 16\,GB
                   (2667\,MT/s) \\
  Almacenamiento & SSD de 234\,GB \\
  Gráfica        & Integrada en el procesador \\
  Hipervisor     & KVM/QEMU con libvirt sobre Ubuntu 26.04 \\
  \bottomrule
\end{tabularx}
\caption{Características del equipo anfitrión}\label{tab:host}
\end{table}
\end{nota}

\section{Herramientas utilizadas}
\begin{table}[H]\centering\small
\begin{tabularx}{\linewidth}{l l X}
  \toprule
  \textbf{Herramienta} & \textbf{Versión} & \textbf{Función en el laboratorio} \\
  \midrule
  KVM / libvirt / virt-manager & --      & Hipervisor y redes virtuales del host \\
  OPNsense                     & 26.7    & Router, cortafuegos, NAT y DHCP entre redes \\
  Suricata (en OPNsense)       & --      & IDS de red sobre la interfaz LAN \\
  Wazuh                        & 4.14    & SIEM/XDR: indexer, manager y dashboard \\
  Sysmon                       & 15.x    & Telemetría avanzada de procesos en Windows \\
  Windows 10 Pro               & 22H2    & Víctima Windows (agente Wazuh + Sysmon) \\
  Ubuntu Server                & 26.04   & Servidor Wazuh y víctima Linux \\
  Kali Linux                   & --      & Máquina atacante en red separada \\
  \bottomrule
\end{tabularx}
\caption{Herramientas y versiones del laboratorio}\label{tab:herramientas}
\end{table}

\section{Estructura del documento}
El capítulo~\ref{cap:arquitectura} describe la arquitectura y las decisiones de
diseño. El capítulo~\ref{cap:despliegue} resume el despliegue y los problemas
encontrados. El capítulo~\ref{cap:escenarios}, núcleo del documento, recoge los
escenarios de ataque con su telemetría, detección, investigación y respuesta.
El capítulo~\ref{cap:cobertura} reúne la cobertura de detección y los ajustes
realizados, y el capítulo~\ref{cap:conclusiones} cierra con lo aprendido. Los
anexos contienen la instalación paso a paso (\ref{anx:instalacion}), los
comandos de gestión (\ref{anx:comandos}) y las configuraciones
(\ref{anx:configuraciones}).
